\documentclass[12pt,a4paper]{article}
\usepackage[margin=1in]{geometry}
\usepackage{amsmath,amssymb,amsthm,mathtools}
\usepackage{booktabs,array,enumitem,fontspec}
\usepackage[hidelinks]{hyperref}
\hypersetup{pdftitle={The Null Cluster as Structural Imprint: A Closed Source Task and the Null-Boundary Programme},pdfauthor={Sidong Liu}}
\defaultfontfeatures{Ligatures=TeX}
\IfFontExistsTF{Times New Roman}{\setmainfont{Times New Roman}}{\setmainfont{TeX Gyre Termes}}
\newtheorem{theorem}{Theorem}[section]
\newtheorem{proposition}[theorem]{Proposition}
\newtheorem{corollary}[theorem]{Corollary}
\theoremstyle{definition}
\newtheorem{definition}[theorem]{Definition}
\newtheorem{conjecture}[theorem]{Conjecture}
\newcommand{\E}{\mathbb E}
\newcommand{\R}{\mathbb R}
\newcommand{\clip}{\operatorname{clip}}
\title{\textbf{The Null Cluster as Structural Imprint}\\
\large A Closed Source Task and the Null-Boundary Programme}
\author{Sidong Liu, PhD\\\small iBioStratix Ltd\\\small\texttt{sidongliu@hotmail.com}}
\date{October 2026 working revision\\\small Series record DOI:
\href{https://doi.org/10.5281/zenodo.20283057}{10.5281/zenodo.20283057}}
\begin{document}
\emergencystretch=2em
\raggedbottom
\maketitle

\begin{abstract}
Null boundaries recur in causal propagation, horizon thermality, entropy
bounds, asymptotic memory and modular geometry. A common interpretation
requires more than their appearance in the same list: one must identify
what a null interface reads, what it omits, and which resources make it
competitive with other interfaces.

We give a closed comparison on a scalar-source family
$S=\theta I+\alpha g$ in a supplied Lorentzian laboratory. Equal-gain
directional probes use the same Gaussian pointers and a common nine-bit
readout. Null contraction removes the unknown trace term. At the specified
mean-square target $1.6$ it needs one record, whereas every adaptive scheme
restricted to either of two non-null causal menus needs at least three;
an eight-record scheme suffices. In the growing-range analogue benchmark
the optimal risks differ by a factor of four. The lower bound allows biased
estimators, analogue access and arbitrary computation. Mixing causal types or allowing
directions sufficiently close to null removes the relevant obstruction.
The null record remains blind to a genuine potential contribution to the
source.

This result establishes a task-specific mechanism for null compression,
not a universal null-data theorem. It motivates a broader hypothesis:
null interfaces are economical where causal transport and quotienting
unwanted source components coincide. We develop that hypothesis against
the distinct mechanisms behind the null cluster and identify source
completion, actual detector costs and propagation tails as decisive tests.
\end{abstract}

\section{The question behind the cluster}
\label{sec:question}

Light cones, horizons and radiation at null infinity do different jobs.
Yet each connects an accessible region to information that reaches its
boundary. This recurrence is worth explaining. The proposed \emph{null
cluster} consists of causal reachability, horizon restriction, thermal
response, entropy accounting, infrared memory, geometric modular flow in
special states, and massless propagation. The facts themselves belong to
established relativity and quantum field theory; the question here is
whether their conjunction identifies a useful principle of physical
readout.

The strongest version would say that null data are the cheapest adequate
interface for a specified physical task. That statement needs a target,
a reconstruction law, a detector menu and a loss function. Comparing
characteristic data reconstructed with the field equation against bulk
samples reconstructed without it would reward unequal knowledge, rather
than the causal type of the interface. Counting the dimensions of a
surface also says little about its bandwidth, sampling or required initial
and boundary data.

We therefore begin with a comparison in which all readers know the same
source law and have access to the same elementary meter. The null reader
wins because its coupling rejects a nuisance source component, not because
its competitor is denied the equations. This gives the programme an
explicit positive instance. Its limits matter equally: the component
discarded as a nuisance for one task is physical source information for
another.

Q-RAIF A supplies a separate algebraic propagation construction with a
Lorentzian causal readout under its stated shell and independence
conditions \cite{Liu2026QRAIFA}. We use a Lorentzian laboratory here as an
input. The present experiment does not derive that laboratory, nor does
the propagation construction by itself grant our pointer couplings. Their
connection is the null contraction and its operational consequences. The
source-completeness question belongs to the companion no-view analysis
\cite{Liu2026NoView}.

\section{A common source and a common meter}
\label{sec:interface}

\subsection{The scalar-source family}

Work in a fixed $1+1$ laboratory with $g=\operatorname{diag}(-1,1)$.
Choose the two-parameter stress family
\begin{equation}
 S=\theta I+\alpha g
   =\begin{pmatrix}\theta-\alpha&0\\0&\theta+\alpha\end{pmatrix},
 \qquad (\theta,\alpha)\in\mathcal B=[0,20]\times[-30,0].
 \label{eq:source}
\end{equation}
The target is $\theta=(S_{00}+S_{11})/2$. The trace coefficient $\alpha$
is unknown and must not be supplied to any estimator.

These tensors have a simple matter realization. For a canonical scalar,
\begin{equation}
 \mathcal L=-\tfrac12 g^{\mu\nu}\partial_\mu\phi\partial_\nu\phi-V,
 \qquad
 S_{\mu\nu}=\partial_\mu\phi\partial_\nu\phi
 -g_{\mu\nu}\bigl(\tfrac12(\partial\phi)^2+V\bigr).
 \label{eq:scalar}
\end{equation}
On the observation patch take $\phi=\sqrt{2\theta}\,t$ and the constant
potential $V=-\alpha\geq0$. The scalar equation is satisfied and
\eqref{eq:scalar} gives \eqref{eq:source}; at $\theta=0$ the field is
constant and the potential term remains. This is the usual scalar stress
formula \cite{Wald1984,Hu2018}, not a new matter equation. Preparation of
the patch, its boundary conditions and the background geometry are common
physical inputs. A homogeneous field on an infinite flat background is
not being claimed as a finite isolated self-gravitating apparatus.

In this family the null energy measures the gradient contribution while
discarding the constant potential. Positivity of that potential does not
make it fictitious: the laboratory energy density is
$S_{00}=\theta+V$. Null sufficiency for $\theta$ is consequently weaker
than source completeness.

\subsection{Directional permission}

An elementary probe is a laboratory-Euclidean unit vector $w$:
$w_0^2+w_1^2=1$. Its signal is
\begin{equation}
 q=g(w,w)=w_1^2-w_0^2,\qquad
 s_q=w^\mathsf T S w=\theta+q\alpha .
 \label{eq:signal}
\end{equation}
A null vector $w=(1,1)/\sqrt2$ gives $q=0$. The two non-null menus are
\begin{equation}
 \mathcal Q_{\rm T}=[-1,-q_0],\qquad
 \mathcal Q_{\rm S}=[q_0,1],\qquad q_0=\tfrac13.
 \label{eq:menus}
\end{equation}
The null menu contains $q=0$. Either sign of the null direction gives the
same signal in this diagonal source family.

The Euclidean normalization describes a calibrated detector gain in this
laboratory; it is not Lorentz invariant. A proper-unit timelike or
spacelike vector of rapidity $\eta$, after this normalization, has
$|q|=\operatorname{sech}(2\eta)$. Thus the non-null menus can be realized
with a supplied boost bound
$|\eta|\leq\tfrac12\operatorname{arcosh}3$. A direction reference,
alignment, gain calibration and this actuator permission are part of the
experiment. Rescaling a vector to increase its signal is not a free
operation. The finite rapidity bound will be used in the proof; removing
it changes the answer.

\subsection{Actual readout and the resource boundary}

Each read uses a fresh independent pointer with canonical coordinates
$Q,P$, $[Q,P]=i\hbar$, initially in
\begin{equation}
 \psi(Q)=(2\pi\sigma^2)^{-1/4}
           \exp[-Q^2/(4\sigma^2)],\qquad \sigma=1.
\end{equation}
The supplied interaction pulse implements
\begin{equation}
 U_q=\exp(-i s_q P/\hbar),\qquad
 H_{\rm int}(t)=\chi(t)s_qP,\qquad\int\chi(t)\,dt=1.
 \label{eq:gate}
\end{equation}
Position measurement of the translated pointer gives
\begin{equation}
 Y_j=\theta+q_j\alpha+Z_j,\qquad Z_j\stackrel{\rm iid}{\sim}N(0,1).
 \label{eq:law}
\end{equation}
This is an actual meter coupling for a stable classical source parameter,
or for a supplied commuting register controlling the translation. It does
not grant a nondisturbing measurement of arbitrary quantum stress. The
front end that transduces \eqref{eq:signal} into the pulse, the pulse clock
and the drive are supplied identically to all menus.

The finite output is a common digitizer. Clip $Y$ to $[-64,64]$, then round
to the nearest multiple of $\Delta=1/2$. There are $257$ possible labels,
stored in nine bits. Equivalently, this is a finite position-bin POVM,
with the two endpoint bins extending to infinity; no intermediate
infinite-precision record is required. The estimator receives only these
labels. Retaining a spent pointer does not grant an additional unpriced
measurement of it. The lower bound below gives its competitors the
stronger permission of full analogue outcomes, so digitization cannot
invalidate it.

All source means obey $-30\leq s_q\leq50$. Thus a common bounded
per-read energy allowance can be assigned to every menu. For example, in
dimensionless oscillator units $H_P=(P^2+Q^2)/2$ and $\hbar=1$, the initial
mean energy is $5/8$ and the translated mean is $5/8+s_q^2/2\leq1250+5/8$.
This allowance is a shared cap, not an assertion that the realized drive
work is the same for every parameter and direction. Preparing the meters,
converting their outcomes, communicating the labels and retaining spent
pointers all cost resources. No reset or autonomous work store is hidden
in \eqref{eq:gate}.

For the comparison we price the number of reads under this fixed
per-read interface, with a maximum pathwise stock of $N$ pointers and
$9N$ output bits. Computation may be arbitrary; it is granted to every
reader, not counted as a null advantage. Choices of $q_j$ may depend on
previous actual records and on independent random seeds. Initial registers
and fresh noise carry no additional information about $\theta,\alpha$.
These conditions exclude an unpriced source oracle and allow fully
adaptive competitors.

\section{Sufficiency and a fair lower bound}
\label{sec:comparison}

\begin{definition}[Closed task]
The task is to estimate $\theta$ in \eqref{eq:source} with worst-case mean
squared error at most $\varepsilon_*=1.6$. A protocol is task-sufficient if
it meets that bound uniformly on $\mathcal B$. Its read cost is the
maximum number of elementary reads on any branch. Let $C_{\rm N}$,
$C_{\rm T}$ and $C_{\rm S}$ be the minimum costs in the three menus.
\end{definition}

This is noisy task sufficiency, rather than exact recovery of the source
tensor. In the analogue null experiment, the sample mean is a sufficient
statistic for the normal location parameter $\theta$. Our actual finite
labels lose some of that information but still meet the task.

\begin{theorem}[Adaptive single-type lower bound]
\label{thm:lower}
For either non-null menu \eqref{eq:menus}, every protocol with at most $N$
reads, including analogue readout, arbitrary randomization and biased
estimation, has worst-case risk at least
\begin{equation}
 \sup_{\mathcal B}\E_{\theta,\alpha}(\widehat\theta-\theta)^2
 \ \geq\ \frac{1}{4\pi^2/20^2+N/4}.
 \label{eq:lower}
\end{equation}
\end{theorem}

\begin{proof}
Restrict the spacelike source family to $\alpha=-3\theta/2$, and the
timelike family to $\alpha=-30+3\theta/2$, with $0<\theta<20$. Both
subfamilies lie inside $\mathcal B$. Revealing the subfamily to the reader
can only help it. Its conditional mean is respectively
\begin{equation}
 \mu_j(\theta)=\theta(1-3q_j/2),\qquad
 \mu_j(\theta)=-30q_j+\theta(1+3q_j/2).
\end{equation}
For either allowed menu the derivative $a_j=\partial_\theta\mu_j$ lies
in $[-1/2,1/2]$.

Fix the independent random seed. At a given history the selection rule
for $q_j$ has no additional dependence on $\theta$. Hence the conditional
Gaussian score is $a_j(Y_j-\mu_j)$ even when the rule is adaptive. These
scores have conditional mean zero and conditional second moment
$a_j^2\leq1/4$. Their cross terms vanish by conditioning on the earlier
history. The score $\ell'$ of the complete transcript therefore obeys
$\E_\theta(\ell')^2\leq N/4$. A stopped protocol can be padded to $N$
reads that the estimator ignores; averaging over the independent seed
preserves the bound.

Use the prior
\begin{equation}
 \pi(\theta)=\frac2{20}\sin^2\frac{\pi\theta}{20},
 \qquad I_\pi=\int_0^{20}\frac{(\pi')^2}{\pi}
              =\frac{4\pi^2}{20^2}.
 \label{eq:prior}
\end{equation}
Clip the estimator to $[0,20]$, which cannot increase its risk. The prior
vanishes at both endpoints. Integration by parts in the joint
prior--transcript density gives
\begin{equation}
 \E\bigl[(\widehat\theta-\theta)(\pi'/\pi+\ell')\bigr]=1.
\end{equation}
The clipped error is bounded, the prior Fisher information is finite,
and the Gaussian score has the square-integrable bound above; the
integration and boundary terms are consequently justified. The
prior--likelihood cross term is zero because every likelihood score has
mean zero. Cauchy--Schwarz yields a Bayes risk of at least
$(I_\pi+N/4)^{-1}$, and worst-case risk dominates Bayes risk.
This is the van Trees argument \cite{GillLevit1995,GassiatPollardStoltz},
applied to the permitted adaptive directional experiment.
\end{proof}

The lower bound does not compare null data with naive volume sampling.
Both non-null readers know the two-parameter source law, can cancel its
trace by combining measurements, and may allocate their reads after
seeing earlier results. Their disadvantage is the restricted gain for
separating $\theta$ from $\alpha$ in one causal type.

\begin{proposition}[Finite-record separation]
\label{prop:finite}
For the common nine-bit readout,
\begin{equation}
 C_{\rm N}=1,\qquad 3\leq C_{\rm T},C_{\rm S}\leq8.
 \label{eq:cost}
\end{equation}
With four reads the null protocol has worst-case risk less than $0.563$,
whereas either non-null menu has risk at least $0.9101698$.
\end{proposition}

\begin{proof}
Let $\varphi(a)$ and $\overline\Phi(a)$ be the standard normal density
and upper tail, and put
\begin{equation}
 B(a)=2\{a\varphi(a)+\overline\Phi(a)\}.
 \label{eq:tail}
\end{equation}
For a signal of absolute value at most $64-a$, clipping contributes an
$L^2$ error at most $\sqrt{B(a)}$: its squared error is bounded by
$Z^2\mathbf1_{\{|Z|>a\}}$, whose expectation is \eqref{eq:tail}.
Rounding adds at most $\Delta/2=1/4$. Thus the per-record digitization
error has $L^2$ norm at most $1/4+\sqrt{B(a)}$.

At $q=0$, $a=44$ works for every source in $\mathcal B$. The clipped
average of $N$ digitized null records therefore has risk at most
\begin{equation}
 \left(N^{-1/2}+\tfrac14+\sqrt{B(44)}\right)^2.
 \label{eq:nullupper}
\end{equation}
This follows by Minkowski from the undigitized sample mean; it requires
no assumption that rounding errors are unbiased or independent. For
$N=1$ the bound is less than $1.563<1.6$; for $N=4$ it is less than
$0.563$. Without any record the minimax risk on $\theta\in[0,20]$ is
$100$, attained by the constant estimate $10$. A random estimate cannot
reduce this value. Hence $C_{\rm N}=1$.

For $N=2$, \eqref{eq:lower} is $1.6702966555\ldots>1.6$; for $N=4$
it is $0.9101698376\ldots$. To prove finite non-null sufficiency, use six
reads at $q=1/3$ and two at $q=1$ for the spacelike menu, or six at
$q=-1/3$ and two at $q=-1$ for the timelike menu. In both cases
\begin{equation}
 \widehat\theta_{\rm an}=\tfrac32\overline Y_{\rm low}
                         -\tfrac12\overline Y_{\rm high}
 \label{eq:endpoint}
\end{equation}
cancels $\alpha$ and has variance $1/2$. All these signals have clipping
margin at least $14$. Replace each average by its digitized average and
clip the final estimate to $[0,20]$. Its risk is at most
\begin{equation}
 \left(\sqrt{\tfrac12}+\tfrac12+2\sqrt{B(14)}\right)^2
 <1.458<1.6.
 \label{eq:nonnullupper}
\end{equation}
Here $B(14)=3.08608544\ldots\times10^{-42}$ and
$B(44)=1.40785656\ldots\times10^{-419}$. This completes the comparison.
\end{proof}

The one-record datum is minimal under deletion of a complete read: its
only proper sub-datum in that menu is no record. Nine bits have not been
proved to be the minimal code length. Further coarsening may still meet
the error target. Likewise, \eqref{eq:cost} bounds the non-null optimum
rather than identifying whether three through seven reads suffice.

The read-cost ratio is now well defined:
\begin{equation}
 \rho_{\varepsilon_*}
 =\frac{C_{\rm N}}{\min(C_{\rm T},C_{\rm S})},
 \qquad \frac18\leq\rho_{\varepsilon_*}\leq\frac13<1.
\end{equation}
It prices actual uses of the common interface. It is not a ratio of
total apparatus energies, preparation costs or optimal detector work.

\section{What produces the advantage, and what removes it}
\label{sec:boundaries}

The causal distinction acts through a concrete statistical geometry.
Null contraction annihilates the entire line of tensors $\alpha g$.
Within either finite-boost non-null menu, the two coefficients are only
partly separable at fixed gain. This explanation is stronger than calling
the boundary thin, but also more conditional: a detector menu determines
which alternatives the word ``cheapest'' covers.

\subsection{The analogue benchmark}

There is a sharp benchmark if the source range is allowed to grow and the
readout is analogue. Let $\theta\geq0$, $\alpha\leq0$ and
$|q|\in[q_0,1]$ within one causal type, with $0<q_0<1$. Replace the
compact prior interval by $[0,L]$. The least-favourable families are
$\alpha=-2\theta/(1+q_0)$ for spacelike directions and
$\alpha=-2L/(1+q_0)+2\theta/(1+q_0)$ for timelike directions. Their gains
are bounded by $r=(1-q_0)/(1+q_0)$. Sending $L$ to infinity gives
\begin{equation}
 \inf_{\rm protocols}\sup_{\theta\geq0,\alpha\leq0}
 \E(\widehat\theta-\theta)^2
 \geq\frac{(1+q_0)^2}{(1-q_0)^2}\frac{\sigma^2}{N}.
 \label{eq:ideal}
\end{equation}
For $q_0=1/3$ and $N$ divisible by four, allocate $3N/4$ reads to the
small-magnitude endpoint and $N/4$ to the other. Formula
\eqref{eq:endpoint} attains $4\sigma^2/N$, so the lower bound is exact
in that analogue benchmark. Null averaging attains $\sigma^2/N$, with
the matching location-model lower bound obtained from the same priors.
The growing-box argument is an ideal statistical benchmark, not an
additional claim about a bounded-energy laboratory.

\subsection{Mixed types and near-null directions}

Allowing both $q=1$ and $q=-1$ changes the permission. Equal allocation
and the mean of the two group averages cancel $\alpha$ and give variance
$\sigma^2/N$, equal to null averaging for even $N$. A physical front end
that measures the sum of these contractions directly is another change
of interface. Neither is included in the single-type lower bound.

The boost restriction is also essential. A one-read non-null estimate
with small $|q|$ has, on $\mathcal B$, digitized root-mean-square error
bounded by
\begin{equation}
 1+\tfrac14+30|q|+\sqrt{B(14)}.
\end{equation}
For $|q|\leq0.0004$ its squared bound is below $1.6$. Thus sufficiently
near-null non-null probes meet the same one-record task. There is no
discontinuous universal advantage attached merely to the sign of
$g(w,w)$. The finite separation is between specified implementable
menus, while nullity gives exact nuisance rejection.

\subsection{The source that the null record omits}

At fixed $\theta$, every value of $\alpha$ gives exactly the same null
output law, for any number of reads. Consider instead estimating the
energy density $S_{00}=\theta-\alpha$. The sources $\alpha=0$ and
$\alpha=-30$ differ in energy density by $30$ but are indistinguishable
to this interface. No nonlinear processing, additional null samples or
more precise null digitizer separates them. Their common output forces
worst-case absolute energy error at least $15$, and worst-case squared
error at least $225$ for those two sources.

This is a blind direction, not a gravitational on/off rule. A detector
that rejects a trace contribution does not remove that contribution from
the matter action. To obtain a complete source certificate one must add
a trace-sensitive measurement, calibrate the missing term from independent
data, or bound it sufficiently for the response task. The companion
analysis \cite{Liu2026NoView} asks exactly when a readout supports such a
certificate. The two responsibilities fit together: economical task
readout and complete physical source accounting need not select the same
observable algebra.

\section{Reading the established null cluster}
\label{sec:cluster}

The source task isolates one mechanism. The wider cluster contains others,
whose distinct assumptions must remain visible. Their conjunction suggests
a research direction rather than a common derivation already in hand.

\subsection{Causal transport and characteristic data}

A relativistic null cone bounds causal influence; it does not confine
every Green function to the cone. Characteristic reconstruction needs
the field equation, gauge constraints and suitable supplementary data
\cite{Wald1984}. Massive fields and curved backgrounds can have tails.
This mechanism removes redundant acquisition using a propagation law,
whereas \eqref{eq:signal} rejects a nuisance at the detector. Combining
them requires one detector and reconstruction analysis.

\subsection{Restricted access, thermality and entropy}

Hawking and Unruh thermality depend on the state, geometry and observation
regime \cite{Hawking1975,Unruh1976}. The vacuum wedge has geometric modular
flow in the Bisognano--Wichmann setting \cite{BisognanoWichmann1975,BisognanoWichmann1976};
an arbitrary restriction need not \cite{Haag1996}. Horizon entropy,
light-sheet bounds and quantum null energy inequalities relate particular
causal boundaries to information \cite{Bekenstein1973,Bousso1999,QNEC2016}.
Quantum focusing remains a conjecture outside its proved regimes
\cite{QFC2016}. Neither these results nor Jacobson's thermodynamic argument
\cite{Jacobson1995} identify finite detector cost with area. Global
entropy, restricted-state entropy and acquired information remain
different quantities.

\subsection{Asymptotic memory and symmetry}

Bondi--Sachs radiation, soft relations and memory connect null-infinity
data to asymptotic charges and persistent detector changes
\cite{Bondi1962,Sachs1962,Weinberg1965,Strominger2018}. Radiative records
can retain a travelling source economically, while constraint, Coulombic
or trace data may still be needed for a complete source certificate.
A retained displacement is an imprint of the passage, not by itself
an observer's self-description.

\subsection{Massless propagation and the finite-resolution edge}

Null momentum characterizes relativistic massless dispersion; gauge
structure can protect masslessness in a specified theory, but does not
follow from nullness alone. The mass--energy companion distinguishes
spectral mass from mass inferred through a readout dictionary
\cite{Liu2026MassEnergy}. Finite resolution near the cone and averaging
momenta rather than invariant masses can change that reading. Particle
masses and mixing still require their dynamical model.

\section{A stronger programme and its tests}
\label{sec:programme}

The finite experiment supplies an answer to a restricted but substantive
question: given this source, this target and these actual detector menus,
null data are sufficient and have a strict read-count advantage. The
larger explanatory claim is that related efficiencies help account for
the recurrence of null structure in physical descriptions. It can be
formulated without making null data universally complete.

\begin{conjecture}[Task-adapted null interfaces]
For important radiation and causal-response tasks, there are physically
implementable null interfaces that attain lower registered resource cost
than the admissible non-null interfaces at the same task error, because
they combine characteristic transport with rejection of source components
irrelevant to that task. The cost advantage survives a common accounting
of detector preparation, alignment, bandwidth, clocks, computation and
any source-completion data required by the target.
\end{conjecture}

The conjecture has two mechanisms and two ways to fail. Characteristic
transport can reduce redundant acquisition, but tails or unmeasured
constraint data can make the chosen boundary inadequate. Nuisance
rejection can improve conditioning, but a mixed or near-null competitor
can match it. Adding realistic direction-setting and meter-renewal costs
could also overturn a sample advantage. These are quantitative competing
explanations, not exceptions to be excluded after a comparison.

A next experiment should use one field preparation, one causal response
target and a common menu of finite surface detectors. Every reconstruction
would receive the same propagation law. The task could be a specified
outgoing waveform or a smeared response, rather than recovery of all bulk
degrees of freedom. The null and non-null reconstructions would be judged
by the same loss and the same resource vector. Their advantage would then
be attributed separately to propagation redundancy and to coupling
conditioning. An interacting background or a field with tails would test
whether source-completion data erase the saving. This extends the present
mechanism rather than promoting it to a theorem about all hypersurfaces.

An interface can be sufficient for one action and incomplete for another.
The null record estimates $\theta$ efficiently; its blindness to $\alpha$
matters when the task becomes total source energy. Thus causal records,
accessible correlations and complete physical response need not select
the same data. Q-RAIF's No View analysis makes this source distinction
explicit \cite{Liu2026NoView}. A trace term unseen by a null instrument
is not thereby removed from gravity. Applications to vacuum energy,
black-hole information or dark sectors require their own source and
retrieval predictions, beyond \eqref{eq:cost}.

\section{Conclusion}

Null contraction can make a physical task cheaper by discarding a source
component that the task does not need. We exhibited the effect with the
same scalar-source family, the same actual pointer law and the same
finite digitizer for all competitors. The proof covers adaptive non-null
readers and records both a sufficient null protocol and finite sufficient
non-null protocols. Its advantage depends on the declared direction
permission and is lost in relevant relaxed menus.

The omitted trace contribution supplies the counterweight: an economical
null interface is not automatically a complete interface to matter or
gravity. The wider null cluster is most promising where causal transport,
restricted access and task-relevant source reduction reinforce one
another. Establishing that conjunction in a common field experiment,
including its completion and renewal costs, is the next test of the
programme.

\begin{thebibliography}{99}
\raggedright
\bibitem{Wald1984} R.~M.~Wald, \emph{General Relativity}, University of Chicago Press (1984).
\bibitem{Hu2018} W.~Hu, \emph{Astro 408: Cosmological Perturbation Theory}, lecture set 2, scalar-field slides, University of Chicago (2018),
\url{https://background.uchicago.edu/~whu/Courses/Ast408_18/ast4PT_02.pdf}.
\bibitem{GillLevit1995} R.~D.~Gill and B.~Y.~Levit, Applications of the van Trees inequality: a Bayesian Cram\'er--Rao bound,
\emph{Bernoulli} \textbf{1}(1/2), 59--79 (1995).
\bibitem{GassiatPollardStoltz} E.~Gassiat, D.~Pollard and G.~Stoltz,
\emph{Revisiting the van Trees inequality in the spirit of H\'ajek and Le Cam}, research draft, Sections 4 and 7.2,
\url{https://stoltz.perso.math.cnrs.fr/Publications/vanTrees-LeCam-TempWeb.pdf}.
\bibitem{Hawking1975} S.~W.~Hawking, Particle creation by black holes,
\emph{Communications in Mathematical Physics} \textbf{43}, 199--220 (1975).
\bibitem{Unruh1976} W.~G.~Unruh, Notes on black-hole evaporation,
\emph{Physical Review D} \textbf{14}, 870--892 (1976).
\bibitem{Bekenstein1973} J.~D.~Bekenstein, Black holes and entropy,
\emph{Physical Review D} \textbf{7}, 2333--2346 (1973).
\bibitem{Bousso1999} R.~Bousso, A covariant entropy conjecture,
\emph{Journal of High Energy Physics} \textbf{07}, 004 (1999).
\bibitem{QFC2016} R.~Bousso, Z.~Fisher, S.~Leichenauer and A.~C.~Wall,
Quantum focusing conjecture, \emph{Physical Review D} \textbf{93}, 064044 (2016).
\bibitem{QNEC2016} R.~Bousso, Z.~Fisher, J.~Koeller, S.~Leichenauer and A.~C.~Wall,
Proof of the quantum null energy condition, \emph{Physical Review D} \textbf{93}, 024017 (2016).
\bibitem{Jacobson1995} T.~Jacobson, Thermodynamics of spacetime: The Einstein equation of state,
\emph{Physical Review Letters} \textbf{75}, 1260--1263 (1995).
\bibitem{BisognanoWichmann1975} J.~J.~Bisognano and E.~H.~Wichmann,
On the duality condition for a Hermitian scalar field,
\emph{Journal of Mathematical Physics} \textbf{16}, 985--1007 (1975).
\bibitem{Haag1996} R.~Haag, \emph{Local Quantum Physics: Fields, Particles, Algebras}, Springer (1996).
\bibitem{Bondi1962} H.~Bondi, M.~G.~J.~van der Burg and A.~W.~K.~Metzner,
Gravitational waves in general relativity VII,
\emph{Proceedings of the Royal Society A} \textbf{269}, 21--52 (1962).
\bibitem{Sachs1962} R.~K.~Sachs, Gravitational waves in general relativity VIII,
\emph{Proceedings of the Royal Society A} \textbf{270}, 103--126 (1962).
\bibitem{Weinberg1965} S.~Weinberg, Infrared photons and gravitons,
\emph{Physical Review} \textbf{140}, B516--B524 (1965).
\bibitem{Strominger2018} A.~Strominger,
\emph{Lectures on the Infrared Structure of Gravity and Gauge Theory}, Princeton University Press (2018).
\bibitem{Liu2026QRAIFA} S.~Liu,
\emph{Lorentzian Read-out from Local Unitary Propagation: Clifford Relations, Causal Cones and Spectral Mass},
2026; series record DOI: 10.5281/zenodo.18525876.
\bibitem{Liu2026NoView} S.~Liu,
\emph{No View from Nowhere: Read-out Closure and the Vacuum Curvature Test in Q-RAIF A},
2026; series record DOI: 10.5281/zenodo.20284454.
\bibitem{Liu2026MassEnergy} S.~Liu,
\emph{Recovering Mass--Energy Equivalence within Algebraic Lorentzian Read-out: Spectral Mass, Lattice Windows and Finite Records},
2026; series record DOI: 10.5281/zenodo.20280951.
\bibitem{BisognanoWichmann1976}
J.~J.~Bisognano and E.~H.~Wichmann, \emph{On the duality condition for quantum fields}, J.\ Math.\ Phys.\ \textbf{17}, 303--321 (1976), doi:10.1063/1.522898.

\end{thebibliography}
\end{document}
